\documentclass[11pt]{amsart}
\usepackage[margin=1in]{geometry}
\usepackage{amsmath,amssymb,amsthm,mathtools}
\usepackage{amscd}
\usepackage[hidelinks]{hyperref}

\title[Three-Dimensional Causal Growth and Octonionic Gauge]{Three-Dimensional Causal Growth\\
and the Octonionic Gauge Sector:\\
Rigidity, Packaging, and Conditional Forcing}
\author{Steven Ettinger Jr}
\address{Independent Researcher}
\email{steven@disregardfiat.tech}
\date{Preprint v1 --- \today}

\newtheorem{theorem}{Theorem}[section]
\newtheorem{lemma}[theorem]{Lemma}
\newtheorem{proposition}[theorem]{Proposition}
\newtheorem{corollary}[theorem]{Corollary}
\newtheorem{definition}[theorem]{Definition}
\newtheorem{remark}[theorem]{Remark}
\newtheorem{conjecture}[theorem]{Conjecture}

\begin{document}
\begin{abstract}
This manuscript is the mathematical sequel to \cite{hqiv-so8-paper}. There, a discrete growth law, a normalized cumulative phase readout, an explicit phase-lift generator \(\Delta\), and a machine-checked Lie closure \(G_2\cup\{\Delta\}\Rightarrow\mathfrak{so}(8)\) were organized into a single pipeline, with careful scope: a locally finite causal partial order does \emph{not}, by itself, determine an octonion multiplication table.

Here we isolate the \emph{dimensional hinge}: the same HQIV null-shell counting law that models \(3{+}1\) Lorentzian causal growth is a \emph{quadratic} polynomial in the shell index, in contrast to the \emph{cubic} ball-volume scaling familiar from coarse \(3{+}1\) CST bookkeeping. We prove elementary combinatorial identities for that law, explain how the quadratic degree matches transverse spatial dimension \(3\) in the HQIV reading, and record the exact machine-checkable \emph{conditional forcing} theorem already packaged in the HQIV-LEAN library: under explicit hypotheses (shared-manifold rapidity bridge, the curvature channel, reference normalization, and the imported \(\mathfrak{so}(8)\) span witness), the conjunction yields a canonical octonionic-layer gauge carrier.

We then develop---as structured narrative, not as a uniqueness certificate---why octonions and \(\mathrm{Spin}(8)\) are the natural \emph{internal} completion of that \(3\)D spatial backbone: Hurwitz rigidity for normed division algebras, the appearance of the Fano plane as the combinatorial spine of imaginary octonion products, and the stacked \(1\)D/\(2\)D/\(3\)D spin-channel interface with optional associativity-breaking language.

We further record a \emph{no-intermediate-knobs} discipline \emph{inside} the growth stack: causal overlap geometry is taken as logically prior to informational monogamy, while Brodie \cite{brodie2026} supplies the thermodynamic route from that monogamy picture to the imprint normalization \(\alpha=3/5\). Concretely, the monogamy coefficient is not a second continuous fit parameter but the complement \(\gamma:=1-\alpha\) on the unit horizon split (\path{gamma_HQIV} in \path{HQVMetric.lean}), and the lattice-derived imprint ratio built from the same cumulative simplex ledger agrees with \(\alpha\) at every shell (\path{latticeAlphaRatio_eq_alpha}). The pair \((\alpha,\gamma)=(3/5,2/5)\) is then the \emph{direct} consequence of that single-causation bookkeeping together with the HQIV normalization and rational arithmetic---not an independent ``Diophantine coincidence'' layered on top. (Lean proves the split and the coherence; the thermodynamic/lattice-dimension narrative for why \(\alpha=3/5\) in nature is in the companion manuscript and repository comments; deriving the same fractions from an \emph{arbitrary} growth law is left as a program, not a claim here.) The \(\mathfrak{so}(8)\) closure stack also admits an exact rational structure-constant certificate. An \emph{audit graph} of named intermediates is posited so that no off-ledger translation steps interpolate between discrete growth and octonionic gauge data; imaginary octonion directions appear as genuine skew-adjoint operators on the shared \(\mathbb{R}^8\) carrier.

Turning this package into a \emph{single} implication from causal-set axioms alone, with no imported Lie seed, remains conjectural and is stated as such.
\end{abstract}
\maketitle

% Shared reader contract: patch QFT + discrete numerics.
% From papers/*.tex:     % Shared reader contract: patch QFT + discrete numerics.
% From papers/*.tex:     % Shared reader contract: patch QFT + discrete numerics.
% From papers/*.tex:     \input{include/patch_theory_messaging.tex}
% From papers/paper/*.tex: \input{../include/patch_theory_messaging.tex}

\paragraph{Patch QFT and discrete numerics (reader contract).}
HQIV is formulated as a \emph{patch theory}: quantum fields, actions, and physical readouts are attached to discrete patches---shell indices $m\in\mathbb{N}$, finite spacetime index types (e.g.\ $\mathrm{Fin}\,4$), and horizon-limited charts---rather than to a hypothetical fundamental smooth spacetime obtained as a mesh-refinement or ultraviolet limit.
Accordingly, when this text refers to ``QFT,'' it means \emph{patch QFT}: patching rules, finite measurement ledgers, and variational identities on the discrete spine; it is not the claim that the theory is \emph{defined} by recovering ordinary continuum quantum field theory through such a limit.
The numbers that admit the sharpest statements---mass and coupling ladders, curvature ratios, shell thermodynamics, and rational witnesses in \texttt{HQIV\_LEAN}---are objects of \emph{discrete mathematics} on $\mathbb{N}$ (combinatorial growth, cumulative simplex ledgers) together with the ordered field $\mathbb{R}$ as used in the proof assistant for analysis, bounds, and phenomenological display.
Smooth-manifold or continuum field notation, when it appears, is a \emph{translation layer} for comparison with standard physics literature; it does not supersede the discrete definitions.

% From papers/paper/*.tex: % Shared reader contract: patch QFT + discrete numerics.
% From papers/*.tex:     \input{include/patch_theory_messaging.tex}
% From papers/paper/*.tex: \input{../include/patch_theory_messaging.tex}

\paragraph{Patch QFT and discrete numerics (reader contract).}
HQIV is formulated as a \emph{patch theory}: quantum fields, actions, and physical readouts are attached to discrete patches---shell indices $m\in\mathbb{N}$, finite spacetime index types (e.g.\ $\mathrm{Fin}\,4$), and horizon-limited charts---rather than to a hypothetical fundamental smooth spacetime obtained as a mesh-refinement or ultraviolet limit.
Accordingly, when this text refers to ``QFT,'' it means \emph{patch QFT}: patching rules, finite measurement ledgers, and variational identities on the discrete spine; it is not the claim that the theory is \emph{defined} by recovering ordinary continuum quantum field theory through such a limit.
The numbers that admit the sharpest statements---mass and coupling ladders, curvature ratios, shell thermodynamics, and rational witnesses in \texttt{HQIV\_LEAN}---are objects of \emph{discrete mathematics} on $\mathbb{N}$ (combinatorial growth, cumulative simplex ledgers) together with the ordered field $\mathbb{R}$ as used in the proof assistant for analysis, bounds, and phenomenological display.
Smooth-manifold or continuum field notation, when it appears, is a \emph{translation layer} for comparison with standard physics literature; it does not supersede the discrete definitions.


\paragraph{Patch QFT and discrete numerics (reader contract).}
HQIV is formulated as a \emph{patch theory}: quantum fields, actions, and physical readouts are attached to discrete patches---shell indices $m\in\mathbb{N}$, finite spacetime index types (e.g.\ $\mathrm{Fin}\,4$), and horizon-limited charts---rather than to a hypothetical fundamental smooth spacetime obtained as a mesh-refinement or ultraviolet limit.
Accordingly, when this text refers to ``QFT,'' it means \emph{patch QFT}: patching rules, finite measurement ledgers, and variational identities on the discrete spine; it is not the claim that the theory is \emph{defined} by recovering ordinary continuum quantum field theory through such a limit.
The numbers that admit the sharpest statements---mass and coupling ladders, curvature ratios, shell thermodynamics, and rational witnesses in \texttt{HQIV\_LEAN}---are objects of \emph{discrete mathematics} on $\mathbb{N}$ (combinatorial growth, cumulative simplex ledgers) together with the ordered field $\mathbb{R}$ as used in the proof assistant for analysis, bounds, and phenomenological display.
Smooth-manifold or continuum field notation, when it appears, is a \emph{translation layer} for comparison with standard physics literature; it does not supersede the discrete definitions.

% From papers/paper/*.tex: % Shared reader contract: patch QFT + discrete numerics.
% From papers/*.tex:     % Shared reader contract: patch QFT + discrete numerics.
% From papers/*.tex:     \input{include/patch_theory_messaging.tex}
% From papers/paper/*.tex: \input{../include/patch_theory_messaging.tex}

\paragraph{Patch QFT and discrete numerics (reader contract).}
HQIV is formulated as a \emph{patch theory}: quantum fields, actions, and physical readouts are attached to discrete patches---shell indices $m\in\mathbb{N}$, finite spacetime index types (e.g.\ $\mathrm{Fin}\,4$), and horizon-limited charts---rather than to a hypothetical fundamental smooth spacetime obtained as a mesh-refinement or ultraviolet limit.
Accordingly, when this text refers to ``QFT,'' it means \emph{patch QFT}: patching rules, finite measurement ledgers, and variational identities on the discrete spine; it is not the claim that the theory is \emph{defined} by recovering ordinary continuum quantum field theory through such a limit.
The numbers that admit the sharpest statements---mass and coupling ladders, curvature ratios, shell thermodynamics, and rational witnesses in \texttt{HQIV\_LEAN}---are objects of \emph{discrete mathematics} on $\mathbb{N}$ (combinatorial growth, cumulative simplex ledgers) together with the ordered field $\mathbb{R}$ as used in the proof assistant for analysis, bounds, and phenomenological display.
Smooth-manifold or continuum field notation, when it appears, is a \emph{translation layer} for comparison with standard physics literature; it does not supersede the discrete definitions.

% From papers/paper/*.tex: % Shared reader contract: patch QFT + discrete numerics.
% From papers/*.tex:     \input{include/patch_theory_messaging.tex}
% From papers/paper/*.tex: \input{../include/patch_theory_messaging.tex}

\paragraph{Patch QFT and discrete numerics (reader contract).}
HQIV is formulated as a \emph{patch theory}: quantum fields, actions, and physical readouts are attached to discrete patches---shell indices $m\in\mathbb{N}$, finite spacetime index types (e.g.\ $\mathrm{Fin}\,4$), and horizon-limited charts---rather than to a hypothetical fundamental smooth spacetime obtained as a mesh-refinement or ultraviolet limit.
Accordingly, when this text refers to ``QFT,'' it means \emph{patch QFT}: patching rules, finite measurement ledgers, and variational identities on the discrete spine; it is not the claim that the theory is \emph{defined} by recovering ordinary continuum quantum field theory through such a limit.
The numbers that admit the sharpest statements---mass and coupling ladders, curvature ratios, shell thermodynamics, and rational witnesses in \texttt{HQIV\_LEAN}---are objects of \emph{discrete mathematics} on $\mathbb{N}$ (combinatorial growth, cumulative simplex ledgers) together with the ordered field $\mathbb{R}$ as used in the proof assistant for analysis, bounds, and phenomenological display.
Smooth-manifold or continuum field notation, when it appears, is a \emph{translation layer} for comparison with standard physics literature; it does not supersede the discrete definitions.


\paragraph{Patch QFT and discrete numerics (reader contract).}
HQIV is formulated as a \emph{patch theory}: quantum fields, actions, and physical readouts are attached to discrete patches---shell indices $m\in\mathbb{N}$, finite spacetime index types (e.g.\ $\mathrm{Fin}\,4$), and horizon-limited charts---rather than to a hypothetical fundamental smooth spacetime obtained as a mesh-refinement or ultraviolet limit.
Accordingly, when this text refers to ``QFT,'' it means \emph{patch QFT}: patching rules, finite measurement ledgers, and variational identities on the discrete spine; it is not the claim that the theory is \emph{defined} by recovering ordinary continuum quantum field theory through such a limit.
The numbers that admit the sharpest statements---mass and coupling ladders, curvature ratios, shell thermodynamics, and rational witnesses in \texttt{HQIV\_LEAN}---are objects of \emph{discrete mathematics} on $\mathbb{N}$ (combinatorial growth, cumulative simplex ledgers) together with the ordered field $\mathbb{R}$ as used in the proof assistant for analysis, bounds, and phenomenological display.
Smooth-manifold or continuum field notation, when it appears, is a \emph{translation layer} for comparison with standard physics literature; it does not supersede the discrete definitions.


\paragraph{Patch QFT and discrete numerics (reader contract).}
HQIV is formulated as a \emph{patch theory}: quantum fields, actions, and physical readouts are attached to discrete patches---shell indices $m\in\mathbb{N}$, finite spacetime index types (e.g.\ $\mathrm{Fin}\,4$), and horizon-limited charts---rather than to a hypothetical fundamental smooth spacetime obtained as a mesh-refinement or ultraviolet limit.
Accordingly, when this text refers to ``QFT,'' it means \emph{patch QFT}: patching rules, finite measurement ledgers, and variational identities on the discrete spine; it is not the claim that the theory is \emph{defined} by recovering ordinary continuum quantum field theory through such a limit.
The numbers that admit the sharpest statements---mass and coupling ladders, curvature ratios, shell thermodynamics, and rational witnesses in \texttt{HQIV\_LEAN}---are objects of \emph{discrete mathematics} on $\mathbb{N}$ (combinatorial growth, cumulative simplex ledgers) together with the ordered field $\mathbb{R}$ as used in the proof assistant for analysis, bounds, and phenomenological display.
Smooth-manifold or continuum field notation, when it appears, is a \emph{translation layer} for comparison with standard physics literature; it does not supersede the discrete definitions.


\section{Introduction and logical scope}

\subsection{What predecessor paper established}

The companion work \cite{hqiv-so8-paper} established a chained template
\[
\text{uniform growth}\Longrightarrow K\Longrightarrow \Omega\Longrightarrow \mathcal{R}
\Longrightarrow \Delta,
\]
together with a machine-checked Lie closure \(\langle G_2,\Delta\rangle_{\mathrm{Lie}}=\mathfrak{so}(8)\) once explicit octonionic matrices are fixed. It also emphasized a methodological point we retain here:
\begin{quote}
\emph{Causal partial order alone does not select the octonion table, the color axis \(e_7\), or the \((e_1,e_7)\) phase plane.}
\end{quote}
Accordingly, any honest use of the word ``forces'' in the present title must be read with a precise hypothesis list.

\subsection{What this paper adds}

We add three layers.

\paragraph{(A) Combinatorial/dimensional rigidity of the HQIV shell law.}
We prove in elementary form that the HQIV shell cardinality
\[
A(m)=4(m+2)(m+1)
\]
has first difference \(\Delta A(m)=8(m+2)\), a \emph{linear} polynomial of degree \(1\) in the shell label \(m\). In the HQIV narrative, this quadratic law is tied to null-cone mode counting in spatial dimension \(3\); bulk ball growth in the same ambient dimension would instead produce a leading \emph{cubic} term. The contrast is not decorative: it isolates which discrete growth law is being used when one speaks of ``\(3\)D causal growth'' in the HQIV sense.

\paragraph{(B) Conditional gauge forcing (Lean packaging).}
We state the library theorem \path{causal_rapidity_forces_octonion} as a transparent implication: if one assumes the SAT shared-manifold bridge and the explicit \(\mathfrak{so}(8)\) span witness, then the conjunction of common rapidity, positivity and divergence of the curvature integral, and reference normalization follows. This is \emph{forcing} in the proof-theoretic sense used throughout formal mathematics: a packaged conclusion from explicit hypotheses, not an ontological claim that Nature must pick \(G_2\).

\paragraph{(C) Octonionic gauge as the natural internal completion.}
We develop, with classical references, why the \(8\)-dimensional normed division algebra \(\mathbb{O}\) and its automorphism group \(G_2\) sit at a special junction once a \(3\)-dimensional spatial rotation story is promoted to an internal holonomy sector requiring triality-compatible bookkeeping. This section is interpretive algebra and representation theory, aligned with the HQVM octonion engine, but it is \emph{not} claimed to be derived from poset axioms without imports.

\paragraph{(D) Growth coherence, unit split, and audited intermediates.}
The pair \((3/5,2/5)\) is the \(d=3\) case of a dimension-indexed family \((\alpha_d,\gamma_d)=\bigl(d/(2d-1),(d-1)/(2d-1)\bigr)\) forced by unit split together with the balance \(\alpha_d/\gamma_d=d/(d-1)\) (Propositions~\ref{prop:balance-spanning-tree} and~\ref{prop:dim-split-family}); the \(d=1\) and \(d=2\) rows are \((1,0)\) and \((2/3,1/3)\). The Lean library then packages the \(d=3\) row: no separate \(\gamma\) knob (\(\gamma=1-\alpha\)), lattice--imprint coherence (\path{latticeAlphaRatio_eq_alpha}), and numeric agreement with \path{alpha_eq_3_5} / \path{alpha_eq_lattice_rational}. Exact rational Lie structure constants sit alongside the floating-point closure table. Methodologically, only book-kept morphisms count as physical intermediates; the seven imaginary units are left-multiplication matrices generating the \(G_2\) sector on \(\mathbb{R}^8\), not an auxiliary \(\mathbb{R}^7\) bolted on after the fact.

\subsection{Reader contract}

\textbf{Proved here (elementary or cited from Lean as indicated):} shell identities; robustness of the curvature channel in \(\alpha>0\); the packaged conditional theorem matching \cite{hqiv-so8-paper} Appendix~A; Propositions~\ref{prop:subadditivity-monogamy}--\ref{prop:balance-spanning-tree} (causal monogamy template and spanning-tree balance, as mathematical arguments); Proposition~\ref{prop:dim-split-family} (algebraic \((\alpha_d,\gamma_d)\) family); Theorems~\ref{thm:growth-coherence} and~\ref{thm:unit-split} as formalized in \path{AlphaGammaForcedByLattice.lean} / \path{OctonionicLightCone.lean} / \path{HQVMetric.lean}.

\textbf{Imported as explicit hypotheses in the formal package:} the \(28\)-dimensional \(\mathfrak{so}(8)\) span witness built from the concrete \(G_2\cup\{\Delta\}\) seed.

\textbf{Classical external theorems used as black boxes in the narrative:} Hurwitz classification of finite-dimensional normed division \(\mathbb{R}\)-algebras; standard facts about \(G_2=\mathrm{Aut}(\mathbb{O})\) and triality for \(\mathrm{Spin}(8)\).

\textbf{Open (stated as conjecture):} uniqueness of the octonionic gauge sector from causal-set axioms and shell asymptotics alone.

\subsection{Executable companion (Lake targets and certificates)}
\label{subsec:lake-targets}
The HQIV-LEAN checkout (\cite{hqiv-lean}) splits heavy certification across named \texttt{lake\ build} targets in \path{lakefile.toml}; this avoids conflating the default formal cone with optional matrix tables.
\begin{itemize}
  \item \textbf{Manuscript / appendix cone:} \texttt{lake\ build\ HQIVPaperClaims} elaborates the growth / curvature / causal-forcing chain together with the \emph{symbolic} \(\mathfrak{so}(8)\) interface (\path{Hqiv/SO8ClosureSymbolic.lean}), without the floating-point Lie-bracket cell cone.
  \item \textbf{Default library cone:} \texttt{lake\ build\ HQIVLEAN} is the main import surface used by narrative modules such as \path{Hqiv/Story/CausalRapidityForcing.lean}; it includes triplet-level \emph{color chart} scaffolding (\path{Hqiv/Physics/StrongColorSu3ChartClosure.lean}, \path{Hqiv/Physics/StrongColorCarrierClosure.lean}) but \emph{excludes} the heavy \(\mathfrak{so}(8)\) matrix-certificate slices and the optional \(\mathfrak{su}(3)\) \(f^{abc}\) simp certificate (see below).
  \item \textbf{Full \(\mathfrak{so}(8)\) matrix closure:} \texttt{lake\ build\ HQIVSO8Closure} rebuilds \path{Hqiv/GeneratorsLieClosureData*.lean}, \path{Hqiv/GeneratorsLieClosure.lean}, and related bracket cells---the setting for Definition~\ref{def:so8-span-witness} once one wants the numeric certificate in-repo.
  \item \textbf{Optional \(\mathfrak{su}(3)\) color certificate:} \texttt{lake\ build\ HQIVStrongColorSu3Certificate} imports the auto-generated \texttt{@[simp]} table for nonzero structure constants \(f^{abc}\) on the active \(3\times 3\) Gell--Mann chart (\path{Hqiv/Physics/StrongColorSu3fStructureSimp.lean}; regenerate via \path{scripts/gen_strong_color_su3_f_simp.py}). The full eight-generator matrix Lie law on all chart pairs remains a documented \emph{target} layered on that table, not a single monolithic proof in the default cone.
  \item \textbf{CI / curated physics surfaces:} continuous integration also builds \texttt{HQIVMeaningfulPhysics}, \texttt{HQIVStory}, and the SU(3) certificate target; see \path{.github/workflows/lean_action_ci.yml}.
\end{itemize}
For a frozen artifact bundle collecting symbolic scripts and JSON excerpts alongside the SO(8) appendix narrative, see \path{papers/so8_closure_full_appendix.tex} (and the discussion of bundle scope in \path{papers/hqiv_rapidity_manifold_so8_closure.tex}).

\section{Quadratic shells versus cubic bulk in \texorpdfstring{$3{+}1$}{3+1} dimensions}
\label{sec:quadratic-vs-cubic}

Fix \(m\in\mathbb{N}\) as a discrete shell index away from a reference event on a locally finite poset, interpreted either abstractly or as an HQIV null-lattice label.

\begin{proposition}[HQIV shell law and increment]
\label{prop:shell-law}
For all \(m\in\mathbb{N}\),
\[
A(m)=4(m+2)(m+1)=4m^2+12m+8,\qquad
A(m+1)-A(m)=8(m+2).
\]
\end{proposition}

\begin{proof}
Expand and subtract.
\end{proof}

\begin{remark}[Degree bookkeeping]
The leading term of \(A(m)\) is \(4m^2\): the shell cardinality is asymptotically quadratic in the index. For a Euclidean ball of radius \(R\) in \(\mathbb{R}^3\), the volume scales as \(R^3\). Thus, if one's discrete model counts \emph{bulk} elements in spatial balls, one expects a cubic leading growth in comparable indices; if one's model counts \emph{null} modes unlocked shell-by-shell on a fixed light-cone in \(3{+}1\) dimensions, the HQIV choice is the quadratic template of Proposition~\ref{prop:shell-law}. The present paper adopts the HQIV null-shell reading throughout.
\end{remark}

\begin{lemma}[Positivity of coefficients in the increment]
For all \(m\in\mathbb{N}\), \(A(m+1)-A(m)>0\).
\end{lemma}

\begin{proof}
\(m+2\ge 2\) and the factor \(8\) is positive.
\end{proof}

\section{Curvature channel and normalized phase readout (summary)}
\label{sec:curvature-summary}

We recall the analytic layer from \cite{hqiv-so8-paper} only as much as needed for self-contained statements.

Fix \(\alpha>0\) (HQIV highlights \(\alpha=\tfrac35\)). Define
\[
\rho(x):=\frac{1}{x}\bigl(1+\alpha\log x\bigr)\quad (x\ge 1),\qquad
K(n):=\sum_{m=0}^{n-1}\rho(m+1).
\]
Then \(K(n)\ge H_n\) (harmonic numbers) and \(K(n)\to+\infty\). After choosing any reference index \(m_\ast\) with \(K(m_\ast)>0\), set \(\Omega(n):=K(n)/K(m_\ast)\). A phase map \(\mathcal{F}\) satisfying the minimal axioms of \cite{hqiv-so8-paper} yields the normalized cumulative phase readout \(\mathcal{R}\).

\begin{theorem}[Uniform-growth curvature theorem --- as in companion paper]
For all \(n\ge 1\), \(K(n)\ge H_n\). Hence \(K(n)\to+\infty\).
\end{theorem}

\begin{proof}
Identical to \cite{hqiv-so8-paper}, Theorem~2.3.
\end{proof}

\section{Conditional forcing: causal growth package \texorpdfstring{$\Rightarrow$}{=>} octonionic-layer witness}
\label{sec:conditional-forcing}

Let \(M\) be a topological space carrying the SAT shared-manifold smooth-bridge data \(B\) used in the Lean development (rapidities, shared profiles, and the clause/variable profile equalities as in \cite{hqiv-so8-paper}, Appendix~A).

\begin{definition}[Imported \(\mathfrak{so}(8)\) span witness]
Let \(\{Y_k\}_{k=0}^{27}\) denote the concrete \(8\times 8\) skew-adjoint generators packaged as \texttt{Hqiv.so8Generator} in HQIV-LEAN. Write \(\mathfrak{s}:=\mathrm{span}_{\mathbb{R}}\{Y_k\}\subset \mathfrak{so}(8)\). The \emph{imported witness hypothesis} is:
\[
\dim_{\mathbb{R}}\mathfrak{s}=28
\quad\text{and}\quad
Y_k\in \mathfrak{s}\ \text{for all }k
\quad(\text{equivalently, each generator lies in the span of the family}).
\]
This hypothesis encodes the already machine-checked closure pipeline \(G_2\cup\{\Delta\}\Rightarrow\mathfrak{so}(8)\) at the level of a finite spanning family.
\end{definition}

\begin{theorem}[Causal rapidity forces octonionic-layer data (Lean)]
\label{thm:lean-forcing}
Assume:
\begin{enumerate}
  \item \(B\) is an \path{SATSharedManifoldSmoothBridge} on \(M\);
  \item the imported \(\mathfrak{so}(8)\) span witness of Definition~4.1 holds for \path{Hqiv.so8Generator}.
\end{enumerate}
Then there exists a shared rapidity observable \(\rho\) on \(M\) matching the variable and clause rapidity profiles of \(B\), and simultaneously the curvature integral is strictly positive at \path{referenceM}, tends to \(+\infty\) along the cofinite filter on \(\mathbb{N}\), satisfies reference-shell normalization \path{omega_k_partial referenceM = 1}, and the \(\mathfrak{so}(8)\) span witness of Definition~4.1 still holds.
\end{theorem}

\begin{proof}
This is exactly the theorem \path{Hqiv.Story.causal_rapidity_forces_octonion} in \path{Hqiv/Story/CausalRapidityForcing.lean}, which composes \path{shared_manifold_induces_common_rapidity}, positivity and divergence of the curvature integral, reference normalization, and the hypothesis \path{hSo8Closure}.
\end{proof}

\begin{remark}[Aliases and narrative names]
The same conjunction is exported under the definitional aliases \path{causal_set_growth_forces_spin8} and \path{causal_growth_forces_real_imaginary_layers_d3}. The discrete increment identity matching Proposition~\ref{prop:shell-law} is \path{causal_shell_growth_law} (alias of \path{new_modes_succ} in \path{OctonionicLightCone.lean}).
\end{remark}

\begin{corollary}[Stacked spin channels with optional superposition interface]
Under the hypotheses of Theorem~\ref{thm:lean-forcing}, and given arbitrary propositions \(P\) (``mild \(3\)D associativity breaking'') and \(Q\) (``superposition-form observable'') together with an implication \(P\Rightarrow Q\), the conjunction of Theorem~\ref{thm:lean-forcing} with \(P\Rightarrow Q\) is also proved as \path{stacked_spin_channel_forcing_d3_with_superposition_interface}.
\end{corollary}

\begin{proof}
Immediate from the proof of that theorem in the same Lean module: it appends \((P\Rightarrow Q)\) to the conclusion without modifying the analytic or Lie data.
\end{proof}

\section{What ``octonionic gauge'' means in this program}

\subsection{Gauge group versus Lie algebra certificate}

Classically, a (compact) gauge group \(G\) in a Euclidean-sector quantum field theory is a Lie group of internal symmetries; infinitesimally, one studies \(\mathfrak{g}=\mathrm{Lie}(G)\). The HQIV-LEAN closure certificate is primarily an \(\mathfrak{so}(8)\) \emph{Lie algebra} statement about explicit matrices, with \(\mathrm{Spin}(8)\) and \(\mathrm{SO}(8)\) used as the corresponding simply connected / adjoint-type Lie groups in the modeling layer.

\begin{definition}[Octonionic gauge sector (HQIV usage)]
The \emph{octonionic gauge sector} is the tuple
\[
(\mathbb{O},\ G_2=\mathrm{Aut}(\mathbb{O}),\ \mathrm{Spin}(8),\ \Delta,\ \mathfrak{so}(8))
\]
in the fixed HQVM/Fano basis conventions: imaginary octonions identified with \(\mathbb{R}^7\), color axis \(e_7\), phase-lift plane \(\mathrm{span}\{e_1,e_7\}\), and \(\Delta\) the corresponding embedded generator in \(\mathfrak{so}(8)\).
\end{definition}

\subsection{From spatial \texorpdfstring{\(\mathrm{SO}(3)\)}{SO(3)} to internal \texorpdfstring{\(\mathrm{Spin}(8)\)}{Spin(8)}}

Spatial rotations in \(3\) dimensions form \(\mathrm{SO}(3)\); quantum mechanically one passes to \(\mathrm{SU}(2)\cong \mathrm{Spin}(3)\). That story is associative and quaternionic. The octonionic gauge sector is \emph{not} the claim that \(\mathrm{SO}(3)\) equals \(\mathrm{Spin}(8)\); rather, it is the claim that a \emph{distinct} internal eight-dimensional carrier, with its own triality automorphisms, is the natural repository for additional internal degrees of freedom once one leaves pure \(3\)D spatial rotations and demands a division-algebra-compatible multiplication law coordinating seven imaginary directions.

\subsection{Hurwitz rigidity (external)}

\begin{theorem}[Hurwitz, classical]
The finite-dimensional normed division algebras over \(\mathbb{R}\) are \(\mathbb{R}\), \(\mathbb{C}\), \(\mathbb{H}\), and \(\mathbb{O}\) only, with real dimensions \(1,2,4,8\).
\end{theorem}

\begin{remark}
Thus, if one's modeling desiderata include: (i) a real normed division algebra larger than the quaternions, and (ii) compatibility with a seven-dimensional imaginary subspace and a natural \(G_2\) symmetry group, then \(\mathbb{O}\) is the unique classical candidate. This is \emph{algebra classification}, not a theorem of causal posets.
\end{remark}

\subsection{Fano combinatorics as the \texorpdfstring{\(3\)D}{3D} projective spine}

The multiplication table of imaginary octonion units is organized by the Fano plane \(\mathrm{PG}(2,2)\): seven points, seven lines, three points per line. This is the same finite geometry that underlies the HQVM labeling of commutator shells. The appearance of a \(2\)-dimensional projective space over \(\mathbb{F}_2\) is a \emph{combinatorial} shadow of ``\(3\)'' (dimension of a vector space over \(\mathbb{F}_2\) whose projectivization is a Fano plane). We record this only as structural poetry linking \(3\)D causal bookkeeping to octonion products, not as a poset-derived theorem.

\section[Growth coherence and audited intermediates]{Growth coherence, audited intermediates, and real work in imaginary octonions}
\label{sec:diophantine-audit}

This section records the \emph{causal order} used in the repository: the \((3/5,2/5)\) numbers are outputs of a growth-internal rule ``one budget, one imprint exponent, one complement,'' not free-floating rationals justified only \emph{a posteriori} by integer identities. Integer combinatorics enters as the \emph{mechanism} that proves the constructed ledger ratio cannot drift from \(\alpha\) shell-by-shell; it is not claimed to be the metaphysical \emph{source} of \(\alpha\) in isolation.

\subsection{Causality forces monogamy (overlap geometry)}

Let \(\mathcal{C}\) be a locally finite causal set with strict partial order \(\prec\). Write \(\operatorname{Past}(e):=\{x\in\mathcal{C}:x\prec e\}\) and \(\operatorname{Fut}(e):=\{x\in\mathcal{C}:e\prec x\}\).

\begin{definition}[Horizon overlap]
\label{def:horizon-overlap}
For \(e_1,e_2\in\mathcal{C}\) define
\[
O(e_1,e_2):=\bigl(\operatorname{Past}(e_1)\cap \operatorname{Past}(e_2)\bigr)\ \cup\ \bigl(\operatorname{Fut}(e_1)\cap \operatorname{Fut}(e_2)\bigr).
\]
\end{definition}

\begin{proposition}[Subadditivity on overlaps (informational monogamy template)]
\label{prop:subadditivity-monogamy}
Let \(I:\mathcal{P}_{\mathrm{fin}}(\mathcal{C})\to \mathbb{R}_{\ge 0}\) be any information functional that is additive on \emph{disjoint} causal domains in the sense that \(I(A\sqcup B)=I(A)+I(B)\) whenever no chain in \(A\) crosses a chain in \(B\) (no shared horizons). If \(O(e_1,e_2)\neq\emptyset\), then necessarily
\begin{equation}
\label{eq:subadd}
I(e_1\cup e_2)\ \le\ I(e_1)+I(e_2)-I\bigl(O(e_1,e_2)\bigr),
\end{equation}
where \(I(e)\) abbreviates \(I(\{x\in\mathcal{C}:x\prec e\}\cup\{e\}\cup\{x:e\prec x\})\) on the chosen causal neighbourhood and the subtracted term is the shared vacuum-correlation budget supported on the overlap.
\end{proposition}

\begin{proof}[Proof sketch]
If \eqref{eq:subadd} failed, the same correlation could be imprinted twice through two distinct chains that later merge, producing either a logical double count of probability mass on overlapping horizons or a putative closed influence loop incompatible with irreflexivity and transitivity of \(\prec\). The inequality is therefore the information-theoretic shadow of acyclicity on overlapping horizons. A full measure-theoretic formalisation is not attempted here.
\end{proof}

\begin{remark}[Lean status]
\label{rem:causal-subadd-lean}
Definition~\ref{def:horizon-overlap} and Proposition~\ref{prop:subadditivity-monogamy} are \emph{templates}: they are not yet specialised to the HQIV readouts \(K\) or \(\mathcal{R}\) as a single packaged Lean theorem from poset axioms alone. The Lean definitions \path{gamma_HQIV := 1 - alpha} and \path{gamma_forced_as_complement} record the downstream \emph{bookkeeping} once a single imprint exponent \(\alpha\) is fixed.
Separately, the triplet-level \(\mathfrak{su}(3)\) chart (\path{Hqiv/Physics/QuarkColorCarrierGaugeScaffold.lean}, \path{StrongColorSu3ChartClosure.lean}, carrier embed \path{StrongColorCarrierClosure.lean}) packages color generators and structure constants \(f^{abc}\) on the active \(3\times 3\) slice of the same octonion-indexed carrier; optional machine-assisted \(f^{abc}\) atoms are under \texttt{HQIVStrongColorSu3Certificate} (subsection~\ref{subsec:lake-targets}).
\end{remark}

\subsection{Derivation of the dimension-balance hypothesis}

Granting the monogamy template of Proposition~\ref{prop:subadditivity-monogamy}, one models how many \emph{independent} overlap interfaces are needed to connect \(d\) null-transverse channels without double-counting shared correlations.

\begin{proposition}[Dimension balance from a minimal overlap graph]
\label{prop:balance-spanning-tree}
Fix an integer \(d\ge 2\). Model the \(d\) independent null-transverse directions as the vertex set \(V=\{v_1,\dots,v_d\}\) of a simple undirected graph \(G=(V,E)\). Interpret each vertex as one independent ``new'' null-transverse channel that can contribute to the curvature-channel imprint (the \(\alpha\log x\) term in \(\rho(x)=\frac1x(1+\alpha\log x)\)). Then the total imprint budget is proportional to \(|V|=d\).

Assume that any consistent assignment of vacuum information across these directions must respect overlaps of causal horizons: in the poset, overlaps appear as shared past/future cones between pairs of directions. The minimal acyclic structure that connects all \(d\) directions without double-counting the same correlation (the monogamy constraint) is a spanning tree \(T\subseteq G\) on \(V\), which has exactly \(|E(T)|=d-1\) edges. Each edge represents one independent monogamy interface linking two directions and consuming one unit of the complementary overlap budget.

Normalising the total horizon budget to \(1\) yields nonnegative weights \(\alpha_d,\gamma_d\) with imprint weight \(\propto d\) and overlap weight \(\propto d-1\), hence
\begin{equation}
\label{eq:balance-system}
\alpha_d+\gamma_d=1,\qquad \frac{\alpha_d}{\gamma_d}=\frac{d}{d-1}.
\end{equation}
This is not an arbitrary Diophantine fit: it is the unique linear system forced by \emph{(i)} unit-split normalisation and \emph{(ii)} the minimal connected overlap pattern on \(d\) transverse directions. For \(d=3\) the tree has \(3\) vertices and \(2\) edges, i.e.\ imprint-to-overlap weight ratio \(3:2\) before normalisation, matching \(\alpha_3:\gamma_3\) in \eqref{eq:balance-system}.
\end{proposition}

\begin{proof}
A tree on \(d\) labelled vertices has exactly \(d-1\) edges (standard graph theory). With imprint mass proportional to \(d\) and overlap mass proportional to \(d-1\), the ratio \(\alpha_d/\gamma_d\) equals \(d/(d-1)\). Imposing \(\alpha_d+\gamma_d=1\) fixes both fractions uniquely.
\end{proof}

\begin{remark}[Why a tree, and not a denser graph]
\label{rem:tree-not-complete}
A cycle or a complete graph would add redundant edges and over-count shared information, violating the ``no double-counting'' clause already enforced by acyclicity of the causal partial order. The spanning-tree pattern is the minimal cycle-free connector compatible with that acyclicity. This matches the ``complementary \((d-1)\)-dimensional overlap weight'' language used elsewhere in the manuscript.
\end{remark}

\subsection[Dimensional derivation]{Dimensional derivation: the same factors for 1D, 2D, and 3D growth}

\begin{proposition}[Imprint pair from unit split and dimension balance]
\label{prop:dim-split-family}
Fix an integer \(d\ge 2\). Suppose nonnegative fractions \(\alpha_d,\gamma_d\) satisfy:
\begin{enumerate}
  \item \textbf{Unit split (no leftover budget):} \(\alpha_d+\gamma_d=1\);
  \item \textbf{Dimension balance:} \(\alpha_d/\gamma_d = d/(d-1)\), as forced by Proposition~\ref{prop:balance-spanning-tree} under the spanning-tree overlap model.
\end{enumerate}
Then
\[
\alpha_d=\frac{d}{2d-1},\qquad \gamma_d=\frac{d-1}{2d-1}.
\]
For \(d=1\) the ratio \(\alpha/\gamma\) is degenerate; the natural boundary reading is \((\alpha_1,\gamma_1)=(1,0)\): a single transverse direction leaves no independent overlap fraction once the unit split is imposed.
\end{proposition}

\begin{proof}
From \(\alpha_d = \frac{d}{d-1}\gamma_d\) and \(\alpha_d+\gamma_d=1\) one gets \(\gamma_d\bigl(\frac{d}{d-1}+1\bigr)=1\), hence \(\gamma_d=\frac{d-1}{2d-1}\) and \(\alpha_d=\frac{d}{2d-1}\).
\end{proof}

\begin{remark}[Low-dimensional specialization]
\label{rem:dim-table}
\[
\begin{array}{c|cc}
d & \alpha_d & \gamma_d\\\hline
1 & 1 & 0\\
2 & 2/3 & 1/3\\
3 & 3/5 & 2/5
\end{array}
\]
Within Proposition~\ref{prop:dim-split-family}, there is at most one consistent pair \((\alpha_d,\gamma_d)\) per integer \(d\ge 2\); the table is therefore exhaustive for that template. The HQIV null lattice in the repository is the \(d=3\) stars-and-bars model \(x+y+z=m\) (\path{latticeSimplexCount} in \path{OctonionicLightCone.lean}), so the physical target row is \((3/5,2/5)\). Formalizing parallel ledgers and hockey-stick ratios for \(d=1\) and \(d=2\) in Lean is natural follow-up work.
\end{remark}

\subsection{Why alpha is three fifths at d equals three (two converging routes)}
\label{subsec:alpha-three-fifths}

\textbf{Algebraic row.} Proposition~\ref{prop:dim-split-family} already fixes the closed form \(\alpha_d=d/(2d-1)\); at \(d=3\) this is \(3/5\), with denominator \(5=2\cdot 3-1\) the automatic output of the linear system in Proposition~\ref{prop:balance-spanning-tree} together with unit normalisation.

\textbf{Lattice route (machine-checked).} Theorem~\ref{thm:growth-coherence} proves \(R(n)=\frac{(n+1)(n+2)(n+3)}{5\cdot \mathrm{cum}(n)}=\alpha\) for every shell index \(n\); with \(\alpha=3/5\) this is the exact rational identity \path{lattice_imprint_ratio_forced_three_fifths}. No external fit is needed: the combinatorial ledger ratio is identical to the curvature imprint exponent at every \(n\).

\textbf{Thermodynamic route (Brodie \cite{brodie2026}).} Applying the Jacobson thermodynamic / entropic-gravity identity on causal horizons, in the package developed by Brodie, to the same null-shell growth law and imposing the monogamy template of Proposition~\ref{prop:subadditivity-monogamy} selects the same normalisation \(\alpha=3/5\) for the \(d=3\) HQIV ladder. This is the ``thermodynamic packaging'' cited in \path{OctonionicLightCone.lean} provenance text.

\textbf{Lean convergence.} The certificates \path{alpha_eq_3_5}, \path{alpha_eq_lattice_rational}, and \path{latticeAlphaRatio_eq_alpha} record numerical agreement inside the audited graph; they do not by themselves replace the causal or graph-theoretic derivations above.

\begin{remark}[Causal priority and Brodie packaging]
\label{rem:causality-monogamy-brodie}
\textbf{Causality first.} Overlap geometry from \(\prec\) is logically prior to monogamy: Proposition~\ref{prop:subadditivity-monogamy} formalises monogamy as subadditivity on \(O(e_1,e_2)\). \textbf{Brodie.} Brodie \cite{brodie2026} supplies the thermodynamic route from that monogamy picture to the imprint coefficient \(\alpha=3/5\) once the HQIV null lattice is fixed.
\end{remark}

\begin{remark}[Why ``\(4\)D'' bookkeeping is non-causal relative to this spine]
\label{rem:four-d-noncausal}
Two distinct statements should not be conflated.

\emph{(i) Combinatorial growth class.} Extending stars-and-bars to \emph{four} nonnegative summands (a putative ``\(d=4\)'' null-transverse ledger) raises the \emph{polynomial degree} of the per-shell mode count: the count at shell \(m\) is \(\binom{m+3}{3}\), whose leading term is cubic in \(m\), whereas the HQIV \(d=3\) ledger \((m+2)(m+1)\) is only quadratic. Cumulative hockey-stick identities then produce a fourth ascending factor \((n+4)\) in the closed-form numerator (parallel to \((n+1)(n+2)(n+3)\) at \(d=3\)). This is the same structural tension already recorded in \cite{hqiv-so8-paper}: bulk ball growth in \(3{+}1\) is cubic in radius, while the HQIV \emph{null-unlock} ladder used for curvature is quadratic in the shell label. Importing a \(d=4\) ledger while pretending to retain the audited \(3+1\) null pipeline therefore inserts an \emph{unreal intermediate}: a new growth law not equivalent to the quadratic HQIV hypothesis.

\emph{(ii) Causal interpretation (thesis, not a Lean theorem).} In sequential-growth / sprinkle language, the operational content of ``causal'' is acyclicity of influence together with a stable notion of ``one step outward on the light cone.'' The audited HQIV story fixes ambient \(3{+}1\) Lorentzian causal structure and uses a \(3\)-dimensional transverse null lattice for shell counts. Treating ``four independent null weights'' as an extra knob without upgrading the ambient causal model does not extend the \emph{same} discrete causal object: it changes both combinatorics and the horizon-overlap pattern on which monogamy is defined. In that \emph{relative} sense---relative to the fixed \(3+1\) quadratic spine---one may say that naive ``\(4\)D'' transverse bookkeeping is \textbf{inherently non-causal}: it is not a coherent extension of the audited growth poset, not a claim that no Lorentzian geometry exists in higher dimension.

For \(d\ge 4\) inside the \emph{same} algebraic template \(\alpha_d=d/(2d-1)\), the fractions continue (e.g.\ \(\alpha_4=4/7\)); the point here is that such rows are \emph{not} compatible with retaining the HQIV \(3+1\) null-shell hypothesis and the no-intermediate audit graph simultaneously.
\end{remark}

\subsection{Unit split inside growth (no second knob)}

\begin{theorem}[Unit horizon split]
\label{thm:unit-split}
Let \(\alpha\) be the HQIV curvature imprint exponent \path{alpha} and define the monogamy coefficient \(\gamma:=1-\alpha\) as \path{gamma_HQIV} in \path{Hqiv/Geometry/HQVMetric.lean}. Then \(\gamma=1-\alpha\) by definition (\path{gamma_forced_as_complement} is definitional), and \(\alpha+\gamma=1\) (\path{alpha_add_gamma}). There is no independently tunable continuous \(\gamma\) in the formalism: the entire ``overlap'' sector is bookkeeping-constrained to consume the complement of the imprint fraction on the unit split.
\end{theorem}

\begin{proof}
By \path{gamma_HQIV}'s definition and \path{alpha_add_gamma}.
\end{proof}

\subsection{Lattice--imprint coherence (no intermediate exponent)}

Write \(\mathrm{cum}(n)\) for \path{cumLatticeSimplexCount} in \path{OctonionicLightCone.lean} and define the combinatorial ratio
\[
R(n):=\frac{(n+1)(n+2)(n+3)}{5\cdot \mathrm{cum}(n)}\in\mathbb{Q}\subset\mathbb{R},
\]
matching Lean's \path{latticeAlphaRatio}.

\begin{theorem}[Growth coherence: ledger ratio equals imprint exponent]
\label{thm:growth-coherence}
For every \(n\in\mathbb{N}\), \(R(n)=\alpha\) as real numbers. Equivalently, \path{latticeAlphaRatio_eq_alpha} holds: the same \(\alpha\) that appears in \(\rho(x)=\frac1x(1+\alpha\log x)\) is the one tracked by the cumulative simplex ledger at \emph{every} shell, with no room for an off-ledger ``effective \(\alpha\)'' floating between combinatorics and curvature.
\end{theorem}

\begin{proof}
This is \path{latticeAlphaRatio_eq_alpha} in \path{OctonionicLightCone.lean}, using the hockey-stick identity \(3\cdot\mathrm{cum}(n)=(n+1)(n+2)(n+3)\) and field arithmetic.
\end{proof}

\begin{corollary}[The \((3/5,2/5)\) pair matches the \(d=3\) dimensional derivation]
\label{cor:three-fifths-two-fifths}
Proposition~\ref{prop:dim-split-family} at \(d=3\) predicts \((\alpha_3,\gamma_3)=(3/5,2/5)\). The repository proves the same values as equalities to the formal constants \path{alpha} and \path{gamma_HQIV}: \path{alpha_eq_3_5}, \path{gamma_forced_two_fifths}, and \(R(n)=3/5\) for all \(n\) by \path{lattice_imprint_ratio_forced_three_fifths} (\path{AlphaGammaForcedByLattice.lean}). The factorization \(\alpha=3/(3+2)\) recorded as \path{alpha_eq_lattice_rational} aligns the abstract fraction \(\alpha_3=d/(2d-1)\) with the stars-and-bars denominator \(5=2d-1\) for \(d=3\), as explained in the docstring above \path{alpha} in \path{OctonionicLightCone.lean}.
\end{corollary}

\begin{proof}
For the algebraic row, substitute \(d=3\) in Proposition~\ref{prop:dim-split-family}. For the Lean equalities, use Theorem~\ref{thm:growth-coherence} with \path{alpha_eq_3_5}, then \path{gamma_forced_two_fifths} and \path{lattice_imprint_ratio_forced_three_fifths}; \path{alpha_eq_lattice_rational} rewrites \(\alpha\) as \(3/(3+2)\).
\end{proof}

\begin{remark}[What Lean proves vs.\ what remains to formalise]
The chain ``no \(\gamma\) knob'' \(\Rightarrow\) ``\(R(n)=\alpha\) for all \(n\)'' \(\Rightarrow\) ``once \(\alpha\) is fixed, \(\gamma\) is fixed'' is fully formal in Lean. Propositions~\ref{prop:subadditivity-monogamy}--\ref{prop:balance-spanning-tree} supply a checkable narrative for the balance ratio \(\alpha_d/\gamma_d=d/(d-1)\); packaging them as a single Lean theorem from causal-set axioms alone is deferred. Proposition~\ref{prop:dim-split-family} is pure algebra once that system is accepted. In the code, \(\alpha=3/5\) is still pinned by the literal definition \path{alpha := 0.60} and \path{norm_num} (\path{alpha_eq_3_5})---an engineering convenience that matches the \(d=3\) row of Proposition~\ref{prop:dim-split-family}. A natural refactor would define \path{alpha} from \(d\) and the closed form, then specialise geometry to \(d=3\). Brodie \cite{brodie2026} and the companion HQIV manuscript supply the thermodynamic path; Remark~\ref{rem:causality-monogamy-brodie} records causal priority. The graph model is essentially a typed enumeration of ``\(d\) vertices, \(d-1\) tree edges''; formalising it in Lean is modest additional effort once overlap is fixed. Extending \path{latticeAlphaRatio}-style coherence to \(d\in\{1,2\}\) would make the lower rows of Remark~\ref{rem:dim-table} machine-checkable, not merely algebraic.
\end{remark}

\begin{remark}[Rational form as consequence, not cause]
Once growth is forced to live on a single audited chain, rationals appear because the ledger is integer-native and the split is algebraic. That is weaker than, and logically distinct from, a Diophantine classification problem: the point is \emph{bookkeeping rigidity inside growth}, not an appeal to arithmetic geometry.
\end{remark}

\subsection{Lie bracket coefficients snap to the rationals}

The machine-checked \(\mathfrak{so}(8)\) closure used in \cite{hqiv-so8-paper} is certified in Lean via explicit generator matrices and floating-point residual checks, but the same bracket relations admit an exact arithmetic proof object: numerator/denominator pairs in \(\mathbb{Q}\) for every structure constant, emitted by \path{scripts/generate_symbolic_lie_closure_certificate.py} as \path{artifacts/so8_symbolic_certificate.json}. At the level of methodology, this is a second snap: after linear algebra over \(\mathbb{Q}\), no bracket coefficient lives in an informal intermediate floating field that could drift off-ledger.

\subsection{Audit graph: no unreal translation steps}

\begin{definition}[Audited pipeline (methodological)]
An \emph{audited intermediate} is a morphism that appears as a named definition, theorem, or reproducible artifact in HQIV-LEAN (or an explicitly linked generator script whose output is checked against Lean). The \emph{audit graph} has vertices \(\{\text{shell combinatorics},\ \rho,K,\Omega,\ \mathcal{R},\ \Delta,\ G_2\text{ commutators},\ \mathfrak{so}(8)\text{ basis}\}\) and edges only along these named constructions.
\end{definition}

\begin{remark}[What counts as ``unreal'' here]
Steps that are standard in informal physics but \emph{not} in the audit graph---for example, complexifying an \(\mathbb{R}^7\) bookkeeping vector without specifying its multiplication, or introducing an ad hoc \(\mathbb{C}^n\) auxiliary space whose dimension is not tied to \(\mathrm{cum}(n)\)---are treated as \emph{unreal} for the purposes of this manuscript: they add degrees of freedom that the discrete story does not budget. The posit is that physical translation in the HQIV program should factor through the audited graph.
\end{remark}

\subsection[Imaginary octonions perform real work]{The imaginary octonions perform real work}

The seven imaginary octonion units are not a post hoc label for ``internal directions.'' In the HQVM construction (mirrored in Lean's generator stack), each imaginary unit \(e_i\) (\(i=1,\dots,7\)) enters through a real left-multiplication matrix \(L(e_i)\in M_8(\mathbb{R})\); the \(G_2\) derivations are built from commutators \([L(e_i),L(e_j)]\), and \(\Delta\) is a fixed skew-adjoint rotation mixing \(e_1\) with \(e_7\) on the same \(\mathbb{R}^8\) carrier. Thus:
\begin{itemize}
  \item imaginary directions contribute \emph{actual} Lie brackets outside the raw linear span of \(G_2\cup\{\Delta\}\) (the \(15<28\) obstruction of \cite{hqiv-so8-paper});
  \item the associator \((x,y,z)\mapsto (xy)z-x(yz)\), which vanishes in \(\mathbb{H}\) but not in \(\mathbb{O}\), is available for \(3\)D-layer associativity-breaking interfaces precisely because those seven directions are realized as noncommuting operators, not as inert labels.
\end{itemize}
In short, ``imaginary'' is a historical octave for \(\mathrm{span}\{e_1,\dots,e_7\}\subset\mathbb{R}^8\); the work performed is \emph{real} linear algebra and Lie closure on that carrier.

\section{Toward a true uniqueness theorem: conjectural program}

\begin{conjecture}[Causal uniqueness of octonionic gauge data]
\label{conj:uniqueness}
Suppose a locally finite causal set \(\mathcal{C}\) arises from a suitable sprinkling into a \(3{+}1\) dimensional globally hyperbolic Lorentzian manifold, and suppose its effective null-shell growth asymptotics match the HQIV quadratic increment law of Proposition~\ref{prop:shell-law}. Then there is a canonical construction of an octonion algebra \(\mathbb{O}(\mathcal{C})\) and a compact internal gauge group \(G(\mathcal{C})\) with \(\mathrm{Lie}(G(\mathcal{C}))\cong \mathfrak{so}(8)\), such that the HQVM \(G_2\cup\{\Delta\}\) seed is functorially determined by \(\mathcal{C}\) up to isomorphism.
\end{conjecture}

\begin{remark}[Status of uniqueness conjecture and Lie seed]
Theorem~\ref{thm:lean-forcing} is a \emph{conditional} packaging statement in the opposite direction: given bridge and seed data, certain analytic conclusions hold. Conjecture~\ref{conj:uniqueness} would supply the missing implication from purely discrete/causal hypotheses to the seed itself.

The conjecture remains open because the current package imports the concrete \(G_2\cup\{\Delta\}\) seed (the \(28\)-dimensional \(\mathfrak{so}(8)\) span witness \path{hSo8Closure}). A full derivation would show that quadratic null-shell asymptotics, local finiteness, and the overlap geometry of Proposition~\ref{prop:subadditivity-monogamy} uniquely determine the Fano-plane multiplication table (or an isomorphic copy) on the seven imaginary directions completing the internal \(\mathbb{R}^8\) carrier. No such functorial reconstruction is presently known; the Lean library correctly treats \path{hSo8Closure} as an explicit hypothesis. The graph-theoretic model of Proposition~\ref{prop:balance-spanning-tree} at least pins the dimension \(d=3\) to the correct row \((\alpha_3,\gamma_3)=(3/5,2/5)\) of the algebraic family, but does not yet select the octonion algebra itself.
\end{remark}

\subsection{Exploration of uniqueness under deformation}
\label{subsec:uniqueness-deformation}

This subsection records an elementary robustness check: Conjecture~\ref{conj:uniqueness} is stated for the \emph{specific} HQIV quadratic law \(A(m)=4(m+1)(m+2)\) tied to isotropic \(3\)-transverse-null stars-and-bars counting in \(3{+}1\) Lorentzian geometry (Proposition~\ref{prop:shell-law}). We stress-test what happens when that law is deformed in two representative ways. Counts for \(m\le 15\) below were verified by a short brute-force Python loop (integer enumeration); the qualitative conclusions do not depend on the cutoff.

\paragraph{Anisotropic null lattice (weighted partition).}
Replace the isotropic constraint \(x+y+z=m\) with a stretched Diophantine constraint, for example
\[
1\cdot x+2\cdot y+3\cdot z=m,\qquad x,y,z\in\mathbb{Z}_{\ge 0},
\]
modelling unequal lattice spacings along three null-transverse directions. Let \(A_{\mathrm{w}}(m)\) denote the number of solutions. Table~\ref{tab:stretch} compares \(A(m)=4(m+1)(m+2)\) to \(A_{\mathrm{w}}(m)\) for small \(m\).

\begin{table}[ht]
\centering
\caption{Isotropic HQIV shell law vs.\ stretched \((1,2,3)\)-weighted counts.}
\label{tab:stretch}
\begin{tabular}{r|rr}
\(m\) & \(4(m+1)(m+2)\) & \(A_{\mathrm{w}}(m)\)\\\hline
0 & 8 & 1\\
1 & 24 & 1\\
2 & 48 & 2\\
3 & 80 & 3\\
4 & 120 & 4\\
5 & 168 & 5\\
6 & 224 & 7\\
7 & 288 & 8\\
8 & 360 & 10
\end{tabular}
\end{table}

The isotropic column is a genuine quadratic polynomial in \(m\). The weighted column is \emph{not} given by any quadratic polynomial: its first differences already fluctuate \((0,1,1,1,1,2,1,2,2,\dots)\), characteristic of quasi-polynomial behaviour in weighted Ehrhart theory. Consequently the clean hockey-stick mechanism underlying Theorem~\ref{thm:growth-coherence} and \path{latticeAlphaRatio_eq_alpha}---constant \(\alpha=3/5\) from a fixed rational ratio at every shell---has no reason to survive: any forced ``effective \(\alpha(m)\)'' would vary with \(m\), or one would introduce an extra continuous stretching parameter, contradicting the no-intermediate-knobs audit discipline of Section~\ref{sec:diophantine-audit}. In short, constant-factor anisotropy pushes the model outside the hypothesis class of Conjecture~\ref{conj:uniqueness}.

\paragraph{Different signature: \((2,2)\) vs.\ \((1,3)\).}
In flat \((2,2)\) split-signature Minkowski space, a product null constraint (after a convenient coordinate choice) leads to \emph{independent} one-dimensional partition factors, e.g.
\[
(x+y=m,\quad u+v=m),\qquad x,y,u,v\in\mathbb{Z}_{\ge 0},
\]
whose solution count per shell is
\[
A_{2+2}(m)=(m+1)^2.
\]
This is still quadratic in \(m\), but with a different leading coefficient and combinatorial origin than \(4(m+1)(m+2)\). It naturally aligns with an effective transverse dimension \(d=2\) rather than \(d=3\): Proposition~\ref{prop:dim-split-family} at \(d=2\) predicts \((\alpha_2,\gamma_2)=(2/3,1/3)\), not \((3/5,2/5)\). At the level of internal symmetry, the Fano-plane spine and seven imaginary octonion units are tied to a \(3\)-dimensional projective geometry over \(\mathbb{F}_2\); a \((2,2)\) product ledger sits lower on the Hurwitz ladder (quaternionic/complex analogues), where the \(\mathfrak{so}(8)\) / \(G_2\cup\{\Delta\}\) closure package of \cite{hqiv-so8-paper} is not the relevant minimal certificate.

\paragraph{Summary.}
Within the audited story, the isotropic \(3\)-transverse-null quadratic law is a rigid hypothesis: it supports constant \(\alpha=3/5\) via lattice coherence, the \(d=3\) row of the imprint family, and---in the conjectural program---the octonionic gauge sector. The deformations above either destroy pure polynomial shell law and constant-ratio coherence (anisotropic weights), or retain quadratic growth but move the effective transverse dimension away from three ((2,2) product law), selecting a different division-algebra level. This is the selection behaviour one expects from a genuine uniqueness theorem: the octonionic package is not universal for arbitrary discrete growth, but tied to the HQIV null-shell equivalence class.

\subsection{Three Natural and one Real (synthesis)}
\label{subsec:three-natural-one-real}

This subsection packages the deeper geometric claim toward which Conjecture~\ref{conj:uniqueness} and the deformation check of Subsection~\ref{subsec:uniqueness-deformation} aim: the audited HQIV data are organised as \textbf{three transverse directions counted in \(\mathbb{N}_0\)} (the ``\(3\) Natural'' part) and \textbf{one ordered timelike axis modeled by \(\mathbb{R}\)} (the ``\(1\) Real'' part), not as an ad hoc \((p,q)\) choice.

\begin{theorem}[Three Natural + one Real under the HQIV shell hypothesis]
\label{thm:three-natural-one-real}
Let \(\mathcal{C}\) be a locally finite causal set with strict partial order \(\prec\). Suppose \(\mathcal{C}\) carries a shell decomposition indexed by \(m\in\mathbb{N}\) whose shell cardinalities agree with the exact HQIV law
\[
A(m)=4(m+1)(m+2)\qquad (m\in\mathbb{N}),
\]
equivalently the stars-and-bars model of Proposition~\ref{prop:shell-law}. Then the following hold.
\begin{enumerate}
  \item\label{it:three-natural} \textbf{(Three Natural.)} Transverse null-mode counting is modeled by three independent nonnegative integers \((x,y,z)\in\mathbb{N}_0^3\) satisfying \(x+y+z=m\); the leading shell growth is quadratic in \(m\) with coefficient matching the \(\binom{m+2}{2}\) lattice. In this sense the discrete transverse geometry is \(3\)-dimensional over the naturals.
  \item\label{it:octonion} \textbf{(Fano / octonions.)} The projective geometry \(\mathrm{PG}(2,2)\) attached to \(\mathbb{F}_2^3\) is the unique Fano plane organising seven imaginary units in a normed division \(\mathbb{R}\)-algebra; by Hurwitz's classification \cite{baez2002,springer2000}, the only non-associative finite-dimensional normed division extension is the octonion algebra \(\mathbb{O}\), with \(\mathrm{Aut}(\mathbb{O})=G_2\) and the \(G_2\cup\{\Delta\}\Rightarrow \mathfrak{so}(8)\) closure package of \cite{hqiv-so8-paper} once explicit matrices are fixed. Dimensions \(d=1,2\) yield only \(\mathbb{R},\mathbb{C},\mathbb{H}\); naive transverse \(d\ge 4\) counting leaves the quadratic HQIV ledger class (Subsection~\ref{subsec:uniqueness-deformation}, Remark~\ref{rem:four-d-noncausal}) and does not support the same audited octonionic carrier.
  \item\label{it:one-real} \textbf{(One Real.)} Along any maximal chain in \(\mathcal{C}\), the restriction of \(\prec\) is a total order; monotone embeddings into \(\mathbb{R}\) supply a real-valued causal height (proper-time readout). The curvature integral \(K(n)\) and normalized phase readout \(\mathcal{R}\) of \cite{hqiv-so8-paper} are built on this same ordered axis after reference normalisation. Proposition~\ref{prop:subadditivity-monogamy} ties monogamy to overlap of past/future cones ordered by \(\prec\).
  \item\label{it:imprint} \textbf{(Imprint consistency.)} Proposition~\ref{prop:dim-split-family} at \(d=3\) gives \(\alpha=3/5\), \(\gamma=2/5\); Theorem~\ref{thm:growth-coherence} identifies the combinatorial ledger ratio with \(\alpha\) at every shell.
\end{enumerate}
In particular, within the hypothesis class of the exact HQIV quadratic shell law and classical division-algebra classification, the only knob-free package supporting the \(28\)-dimensional octonionic gauge spine is the \(3{+}1\) template ``three natural transverse coordinates + one real causal height.''
\end{theorem}

\begin{proof}[Proof sketch]
For \eqref{it:three-natural}, expand \(4(m+1)(m+2)=8\binom{m+2}{2}\) and compare with the stars-and-bars count of solutions to \(x+y+z=m\) scaled by the octonion factor recorded in \path{OctonionicLightCone.lean} (see \path{available_modes_octonion}).

For \eqref{it:octonion}, cite Hurwitz and the uniqueness of \(\mathbb{O}\) among non-associative normed division algebras; identify \(\mathrm{PG}(2,2)\) with lines in \(\mathbb{F}_2^3\) as in the Fano discussion of Section~\ref{sec:diophantine-audit}. The exclusion of other transverse dimensions combines the Hurwitz ladder with the deformation obstructions in Subsection~\ref{subsec:uniqueness-deformation}.

For \eqref{it:one-real}, any maximal chain in a locally finite partial order is order-isomorphic to a countable suborder of \(\mathbb{Q}\) or \(\mathbb{R}\) after completion; the HQIV curvature channel is indexed by \(n\in\mathbb{N}\) along that axis.

For \eqref{it:imprint}, this is Proposition~\ref{prop:dim-split-family} at \(d=3\) together with Theorem~\ref{thm:growth-coherence}.
\end{proof}

\begin{remark}[Causal-set dimension estimators and functoriality]
\label{rem:three-natural-one-real-status}
The step ``quadratic growth \(\Rightarrow\) three transverse directions'' can also be rewritten in continuum or sprinkling language using dimension estimators of Myrheim--Meyer type in causal-set theory \cite{bombelli1987,sorkin2005}; in this manuscript that route is treated as an optional translation layer, not as a foundational requirement. The HQIV ontology is discrete and observable-universe-complete, so no continuum limit is required or desired for the core claims. Likewise, item~\eqref{it:octonion} does \emph{not} claim that \(\mathcal{C}\) functorially constructs \(\mathbb{O}\): Conjecture~\ref{conj:uniqueness} remains the precise target for removing the imported \path{hSo8Closure} witness. Theorem~\ref{thm:three-natural-one-real} is therefore a \textbf{synthesis} of audited discrete inputs + classical algebra + the causal templates of Section~\ref{sec:diophantine-audit}; continuum equations are convenience summaries of that discrete structure, not a replacement for the formal uniqueness proof.
\end{remark}

\section{Discussion}

The companion paper \cite{hqiv-so8-paper} already established that the growth-to-curvature-to-readout chain is structurally robust in \(\alpha>0\) and that the Lie closure is a machine-checked certificate once matrices are fixed. The present sequel sharpens the \emph{dimensional narrative}: \(3\)D null transversality motivates the quadratic HQIV shell law, while Hurwitz and \(G_2\) triality motivate the octonionic gauge sector as the classical unique division-algebra completion beyond the quaternions. Section~\ref{sec:diophantine-audit} derives the balance ratio from a spanning-tree overlap model (Proposition~\ref{prop:balance-spanning-tree}), records a causal monogamy template (Definition~\ref{def:horizon-overlap}, Proposition~\ref{prop:subadditivity-monogamy}), and gives \((3/5,2/5)\) as the \(d=3\) entry of the closed family \(\alpha_d=d/(2d-1)\), \(\gamma_d=(d-1)/(2d-1)\) from unit split + balance (Proposition~\ref{prop:dim-split-family}), then anchors that row to Lean via \(\gamma=1-\alpha\) and \path{latticeAlphaRatio_eq_alpha}. Remarks~\ref{rem:causality-monogamy-brodie}--\ref{rem:four-d-noncausal} articulate Brodie's thermodynamic packaging, the exhaustiveness of the \((1,0)\), \((2/3,1/3)\), \((3/5,2/5)\) ladder under that template, and why naive \(d=4\) transverse counting is incompatible with the quadratic \(3+1\) audited spine. Imaginary octonion directions are realized as operators on \(\mathbb{R}^8\), not as an informal auxiliary complexification.

What remains for a fully causal \emph{derivation} of that sector is Conjecture~\ref{conj:uniqueness}: a functorial reconstruction of \(\mathbb{O}\) and the \(G_2\cup\{\Delta\}\) seed from poset-level axioms, rather than an explicit hypothesis. Subsection~\ref{subsec:uniqueness-deformation} records a reproducible robustness check: anisotropic weights and \((2,2)\) product null counting leave the HQIV hypothesis class and break either constant-\(\alpha\) coherence or the \(d=3\) imprint row. Theorem~\ref{thm:three-natural-one-real} then states the resulting \(3{+}1\) ``Natural + Real'' synthesis under the exact shell law, with Remark~\ref{rem:three-natural-one-real-status} (following that theorem) separating what is Lean-audited from what still imports classical algebra or causal-set folklore.

Appendix~\ref{sec:thermo-derivation} supplies a self-contained thermodynamic packaging of the \((\alpha_d,\gamma_d)\) family (Jacobson horizon identity as read through Brodie \cite{brodie2026}), aligned with the monogamy template of Proposition~\ref{prop:subadditivity-monogamy} and the graph-theoretic unit split of Theorem~\ref{thm:unit-split}.

\appendix
\section[Thermodynamic derivation of the imprint family]{Thermodynamic Derivation of the \texorpdfstring{\((\alpha_d,\gamma_d)\)}{(alpha\_d, gamma\_d)} Family}
\label{sec:thermo-derivation}

This appendix supplies the thermodynamic derivation of the dimension-indexed family \((\alpha_d,\gamma_d)\) using the Jacobson horizon identity as packaged by Brodie \cite{brodie2026}, under the causal monogamy template of Proposition~\ref{prop:subadditivity-monogamy}. The goal is to make the ``two converging routes'' claim fully self-contained.

\subsection{Monogamy on the horizon}

Let \(\mathcal{H}\) be a discrete causal horizon patch in the HQIV null lattice, decomposed into \(d\) transverse null directions (the vertices \(V\) of the overlap graph \(G\)). By Definition~\ref{def:horizon-overlap}, any two directions \(e_i,e_j\) that share a past or future cone contribute an overlap set \(O(e_i,e_j)\). The information functional \(I\) (entropy in nats, or number of independent bits) obeys the subadditivity
\[
I(\mathcal{H})\le\sum_{i=1}^d I(e_i)-\sum_{\text{edges }(i,j)}I(O_{ij}),
\]
where the sum over edges runs over the minimal connecting structure. The minimal acyclic connector is the spanning tree \(T\) with exactly \(d-1\) edges (Proposition~\ref{prop:balance-spanning-tree}).

\subsection{Jacobson thermodynamic identity (discrete form)}

Following Jacobson, the entropy production across \(\mathcal{H}\) is tied to the energy-momentum flux through the horizon:
\[
\delta S=\frac{2\pi}{\hbar}\int_{\mathcal{H}}T_{\mu\nu}\chi^\mu\,d\Sigma^\nu,
\]
where \(\chi\) is the approximate boost Killing field normal to the null generators. In the discrete HQIV setting this becomes an equality between the variation of the information functional \(\delta I\) and the accumulated curvature channel \(\rho(x)=x^{-1}(1+\alpha\log x)\). The term \(\alpha\log x\) is precisely the \emph{imprintable} contribution: the part of \(\delta I\) that is not redundantly locked in monogamous overlaps and therefore contributes to the net curvature integral \(K(n)\).

\subsection{Effective imprint fraction}

The total horizon budget is normalised to unity (unit-split condition of Theorem~\ref{thm:unit-split}). Each of the \(d\) transverse directions contributes one independent ``imprint unit'' (a new null mode that can carry free energy flux). Each of the \(d-1\) monogamy interfaces consumes one unit of correlation budget that \emph{cannot} be imprinted independently (the subtracted \(I(O_{ij})\) term). Consequently the effective imprinted information is
\[
I_{\mathrm{imprint}}=d\cdot u,\qquad I_{\mathrm{overlap}}=(d-1)\cdot u
\]
for a common unit \(u>0\). The total budget is therefore
\[
I_{\mathrm{total}}=I_{\mathrm{imprint}}+I_{\mathrm{overlap}}=(2d-1)u=1\quad\Rightarrow\quad u=\frac{1}{2d-1}.
\]
The imprint coefficient that multiplies the logarithmic term in the curvature channel is therefore the imprinted fraction
\[
\alpha_d=\frac{I_{\mathrm{imprint}}}{I_{\mathrm{total}}}=\frac{d}{2d-1},
\]
with the complementary monogamy coefficient
\[
\gamma_d=1-\alpha_d=\frac{d-1}{2d-1}.
\]
This is exactly the algebraic family already obtained from the pure graph-theoretic balance (Proposition~\ref{prop:dim-split-family}).

\subsection{Convergence of routes and specialisation to \texorpdfstring{\(d=3\)}{d=3}}

The thermodynamic route (Jacobson identity + monogamy subadditivity on the spanning-tree overlap graph) therefore yields \emph{identical} fractions to the lattice-coherence route (hockey-stick identities on the stars-and-bars count). For the physical HQIV case \(d=3\) one recovers the canonical value
\[
\alpha_3=\frac{3}{5},\qquad\gamma_3=\frac{2}{5}
\]
with no free parameters. The same mechanism supplies the lower-dimensional rows \((1,0)\) and \((2/3,1/3)\) and predicts the generalisation to higher transverse dimension (subject to the quadratic-shell constraint of Proposition~\ref{prop:shell-law}).

This closes the loop: the numerical value \(\alpha=3/5\) is not an external fit but the inevitable consequence of (i) causal acyclicity (monogamy), (ii) the Jacobson flux-entropy identity on null horizons, and (iii) the three-dimensional transverse null lattice of the HQIV model.

\section*{Acknowledgments}
AI-assisted drafting and repository tooling were used during manuscript preparation. Formal claims are anchored to the HQIV-LEAN library; the author reviewed the hypothesis lists in \path{Hqiv/Story/CausalRapidityForcing.lean}.

\begin{thebibliography}{99}

\bibitem{hqiv-so8-paper}
S. Ettinger Jr, \emph{From Local Rapidity Fiber to Spin(8): The Primitive Rapidity Observable and 28-Dimensional Closure}, HQIV-LEAN repository, \path{papers/hqiv_rapidity_manifold_so8_closure.tex}, 2026.

\bibitem{hqiv-lean}
HQIV-LEAN repository, \url{https://github.com/disregardfiat/hqiv-lean}, 2026.

\bibitem{baez2002}
J. C. Baez, \emph{The Octonions}, Bull. Amer. Math. Soc. 39 (2002), no.~2, 145--205.

\bibitem{springer2000}
T. A. Springer and F. D. Veldkamp, \emph{Octonions, Jordan Algebras and Exceptional Groups}, Springer, 2000.

\bibitem{bombelli1987}
L. Bombelli, J. Lee, D. Meyer, and R. D. Sorkin, \emph{Space-time as a causal set}, Phys. Rev. Lett. 59 (1987), 521--524.

\bibitem{sorkin2005}
R. D. Sorkin, \emph{Causal sets: Discrete gravity}, in \emph{Lectures on Quantum Gravity}, eds.\ A. Gomberoff and D. Marolf, Springer, 2005, pp.~305--327.

\bibitem{brodie2026}
K. Brodie, \emph{Derivation of Quantised Inertia from Jacobson Thermodynamics}, Zenodo (2026), doi:\,\url{https://doi.org/10.5281/zenodo.18706746}.

\end{thebibliography}

\end{document}
